\subsubsection*{آ: \lr{nginx}}
{کاربرد:} به عنوان \lr{API Gateway}، لود بالانسر و \lr{Reverse Proxy} عمل می‌کند. درخواست‌های ورودی کاربران را دریافت کرده و به میکروسرویس مربوطه هدایت می‌کند. همچنین وظایفی مثل \lr{SSL Termination} و کش کردن محتوای استاتیک را انجام می‌دهد. \\
{مثال واقعی:} روی یک سرور، فرانت‌اند روی \lr{app.example.com} و بک‌اند روی \lr{api.example.com} قرار دارند؛ \lr{nginx} بر اساس ساب‌دامین، درخواست‌ها را به‌ترتیب به سرویس فرانت (مثلا \lr{React} روی پورت ۳۰۰۰) و سرویس \lr{API} (مثلا روی پورت ۸۰۰۰) هدایت می‌کند.

\subsubsection*{ب: \lr{Prometheus}}
{کاربرد:} ابزاری برای مانیتورینگ و جمع‌آوری متریک‌ها به صورت \lr{Pull-based} است. وضعیت سلامتی، میزان مصرف \lr{CPU/RAM} و تعداد درخواست‌های هر میکروسرویس را ثبت می‌کند. \\
{مثال واقعی:} سرویس پرداخت متریکی به نام \lr{successful\_payments\_total} را اکسپوز می‌کند و \lr{Prometheus} هر ۳۰ ثانیه آن را می‌خواند تا در صورت افت ناگهانی تراکنش‌ها، هشدار صادر کند.

\subsubsection*{پ: \lr{RabbitMQ}}
{کاربرد:} یک \lr{Message Broker} مبتنی بر پروتکل \lr{AMQP} است که برای ارتباطات ناهمگام و مطمئن بین سرویس‌ها استفاده می‌شود. پیام‌ها را در صف‌ها نگه می‌دارد تا سرویس مقصد آن‌ها را پردازش کند. \\
{مثال واقعی:} در یک سیستم فروشگاهی، پس از خرید موفق، سرویس سفارش پیامی به صف \lr{shipping-queue} می‌فرستد تا سرویس ارسال کالا بعدا آن را بردارد و پردازش کند.

\subsubsection*{ت: \lr{Celery}}
{کاربرد:} یک صف وظایف توزیع‌شده برای پایتون است که کارهای سنگین، زمان‌بر یا دوره‌ای را خارج از چرخه اصلی درخواست کاربر اجرا می‌کند. \\
{مثال واقعی:} وقتی کاربر درخواست خروجی اکسل از گزارش‌های یک‌ساله خود را می‌دهد، سرویس وب این کار سنگین را به \lr{Celery} می‌سپارد تا در پس‌زمینه انجام شود و به کاربر پیام گزارش در حال آماده‌سازی است را فورا نشان می‌دهد.

\subsubsection*{ث: \lr{ELK (Elasticsearch, Logstash, Kibana)}}
{کاربرد:} پشته‌ای برای مدیریت، تجمیع و تحلیل لوکال و متمرکز لاگ‌ها در سیستم‌های توزیع‌شده است. \lr{Logstash} لاگ‌ها را جمع‌آوری و فیلتر می‌کند، \lr{Elasticsearch} آن‌ها را ایندکس و ذخیره می‌کند و \lr{Kibana} محیط گرافیکی نمایش آن‌هاست. \\
{مثال واقعی:} وقتی یک خطای ناشناخته در فرآیند ثبت‌نام رخ می‌دهد، برنامه‌نویس به جای بررسی لاگ‌های تک‌تک سرورها، در \lr{Kibana} شناسه کاربر را جستجو می‌کند تا تمام لاگ‌های مربوط به او را در میکروسرویس‌های مختلف ببیند.

\pagebreak

\subsubsection*{ج: \lr{Redis}}
{کاربرد:} یک پایگاه داده درون‌حافظه‌ای بسیار سریع است که بیشتر به عنوان لایه کش، مدیریت نشست‌ها یا \lr{Rate Limiter} استفاده می‌شود. \\
{مثال واقعی:} اطلاعات پروفایل کاربران یا کاتالوگ محصولات پربازدید در \lr{Redis} کش می‌شود تا نیازی به کوئری زدن مداوم به دیتابیس‌های اصلی و کندتر نباشد.

\subsubsection*{چ: \lr{Kubernetes}}
{کاربرد:} پلتفرم ارکستراسیون کانتینرها (\lr{Container Orchestration}) است. وظیفه مدیریت چرخه حیات، استقرار اتوماتیک، مقیاس‌پذیری خودکار (\lr{Auto-scaling}) و حفظ پایداری کانتینرهای میکروسرویس‌ها را بر عهده دارد. \\
{مثال واقعی:} در زمان کمپین‌های تخفیفی که ترافیک سایت ۱۰ برابر می‌شود، \lr{Kubernetes} به طور خودکار تعداد \lr{Pod}های سرویس کاتالوگ را افزایش می‌دهد تا سیستم زیر بار ترافیک نخوابد.

\subsubsection*{ح: \lr{Kafka}}
{کاربرد:} یک پلتفرم توزیع‌شده برای پردازش و انتقال جریان داده‌ها (\lr{Event Streaming}) با حجم و سرعت بسیار بالا است. بر خلاف واسط‌های سنتی، پیام‌ها را به صورت پایدار روی دیسک ذخیره می‌کند و چندین مصرف‌کننده می‌توانند تاریخچه داده‌ها را بازخوانی کنند. \\
{مثال واقعی:} در سیستم‌های تاکسی اینترنتی، جریان مختصات جغرافیایی (\lr{GPS}) راننده‌ها به صورت مداوم به \lr{Kafka} فرستاده می‌شود تا سرویس‌های تخمین قیمت، نقشه و بررسی رفتار راننده همزمان از آن استفاده کنند.

\subsubsection*{خ: \lr{Jaeger}}
{کاربرد:} ابزاری برای ردیابی توزیع‌شده (\lr{Distributed Tracing}) است. مسیر حرکت یک درخواست کاربر را در میان میکروسرویس‌های مختلف دنبال کرده و زمان صرف شده در هر مرحله را به صورت گرافیکی نمایش می‌دهد تا گلوگاه‌ها کشف شوند. \\
{مثال واقعی:} اگر لود شدن صفحه سبد خرید ۵ ثانیه طول بکشد، با استفاده از \lr{Jaeger} مشخص می‌شود که مثلا ۴.۲ ثانیه از این زمان مربوط به کندی کوئری دیتابیس در سرویس تخفیف بوده است.


$\\$

\includegraphics[width=1\textwidth]{figs/1.png}